\documentclass[11pt]{article}
\usepackage[margin=1in]{geometry}
\usepackage{amsmath,amssymb,amsthm}
\usepackage{xurl}
\usepackage{hyperref}
\sloppy
\let\oldus\_ \renewcommand{\_}{\oldus\allowbreak}
\newtheorem{theorem}{Theorem}
\newtheorem{lemma}[theorem]{Lemma}
\theoremstyle{remark}
\newtheorem{remark}[theorem]{Remark}
\newcommand{\Q}{\mathbb{Q}}
\newcommand{\Z}{\mathbb{Z}}
\newcommand{\C}{\mathbb{C}}
\newcommand{\OO}{\mathcal{O}}

\title{Conjecture 00000007928 is false:\\ the formula $|K_2(\OO_F)| = |\zeta_F(-1)|\,D_F\,2^{[F:\Q]}$ fails for $F=\Q$}
\author{Jackmeson1}
\date{}

\begin{document}
\maketitle

\section{The conjecture}

The English statement reads:

\begin{quote}\small
Definition: The spectralization of the tame kernel: a Brauer-type explicit formula for the tame
kernel order $w_2(F) = |K_2(\OO_F)|$. Conjecture: The correction to the 2-adic explicit formula
$w_2(F) = \prod_p p^{\sum e(f_p)}$ is the adjustment of Dirichlet units:
$w_2(F) = |\zeta_F(-1)|\, D_F\, 2^{[F:\Q]}$ (an integralization of a rational identity), and the
$p$-part of the formula is given by the valuation sum of the local dilogarithm values of the tame
symbol on units, with the unit contribution obeying a logarithmic law linear in the unit rank of
$F$. (tame kernel Brauer explicit formula)
\end{quote}

The Chinese statement has the same content. It also gives the displayed formula
$w_2(F)=|K_2(\OO_F)|$ and $w_2(F)=|\zeta_F(-1)|\cdot D_F\cdot 2^{[F:\Q]}$, with the same
parenthetical (``integralization of a rational identity'').

\section{Reading}

The Definition line fixes $w_2(F) := |K_2(\OO_F)|$, the order of the group $K_2$ of the ring of
integers $\OO_F$ of a number field $F$. The main clause asserts the equality
\[
  (\ast)\qquad |K_2(\OO_F)| \;=\; |\zeta_F(-1)| \cdot D_F \cdot 2^{[F:\Q]},
\]
for number fields $F$ (the quantifier is implicit; the formula is stated for a general $F$). The
conjecture is a conjunction of $(\ast)$ with further clauses (the $p$-part given by local
dilogarithm values, and a logarithmic law for the unit contribution). We refute $(\ast)$, which
refutes the conjunction. We do not address the other clauses. Our witness is $F=\Q$, a totally
real field whose unit rank is $0$. A reading that restricts $(\ast)$ to fields of positive unit rank
is not refuted here.

\paragraph{Conventions.}
\begin{itemize}
\item $\zeta_F$ is the Dedekind zeta function. It is the Dirichlet series
  $\sum_{n\ge 1} a_n n^{-s}$, where $a_n$ is the number of ideals of $\OO_F$ of absolute norm $n$.
  This series converges for $\operatorname{Re} s>1$. It has a meromorphic continuation to $\C$
  whose only pole is a simple pole at $s=1$ (Wikipedia, ``Dedekind zeta function'', quoted in the
  attempt record). $\zeta_F(-1)$ means the value of this continuation.
\item $D_F$ is the discriminant of $F$, taken in absolute value $|d_F|$. For $F=\Q$ we have
  $d_\Q = 1$, so the choice of sign convention does not matter.
\item $[F:\Q]$ is the degree $\dim_\Q F$; it equals $1$ for $F=\Q$.
\item $|G|$ is the number of elements of the group $G$. In Lean we use \texttt{Nat.card G}, which is
  $0$ for infinite $G$. An infinite group has no finite order, so in particular it does not have
  order $1/6$.
\end{itemize}

\paragraph{The parenthetical ``an integralization of a rational identity''.}
The main clause displays an equality between $w_2(F)$, the order of a group, and the number
$|\zeta_F(-1)|D_F2^{[F:\Q]}$. The parenthetical describes this equality as an integral form of a
rational identity. It does not replace the displayed equality by a different formula, and it names
no integralization procedure. We refute the displayed equality, which is the clause's assertion.
For $F=\Q$ the right-hand side is the non-integer $1/6$ (Theorem~\ref{thm:main}). So $(\ast)$ is not
an integral identity at $F=\Q$, and no integral identity is obtained by asserting it there. A
different reading, in which the left side equals the right side only after some unspecified
operation (for example a numerator or a rounding), is a different formula. It is not refuted here.

\paragraph{Why $K_2$ itself is not built.}
The refutation does not use any property of $K_2(\Z)$. The order of any group is a natural number,
or else the group is infinite. The right-hand side of $(\ast)$ at $F=\Q$ is $1/6$, which is neither.
The Lean theorem therefore quantifies over \emph{every} type $G$ and states that \texttt{Nat.card G}
is not the right-hand side. This covers $G=K_2(\Z)$ whatever its construction. The same argument
covers $w_2$ read in its other usual sense (a positive integer defined through roots of unity), since
that is a natural number too.

\paragraph{Context (not used in the proof).}
The formula $(\ast)$ is not a rewriting of the Birch--Tate conjecture. Wikipedia's Birch--Tate
article, retrieved for this report, says that $K_2$ of a number field ``is also known as the tame
kernel'' and that the conjecture ``relates the order of this group (its number of elements) to the
value of the Dedekind zeta function''. To be precise, the finite group meant here is
$K_2(\OO_F)$, the tame kernel, i.e.\ the kernel of the tame-symbol map from $K_2(F)$ to
$\bigoplus_{\mathfrak p}k(\mathfrak p)^\times$ over the finite primes; it is not $K_2(F)$ itself. The
formula given there, for totally real $F$, is $\#K_2(\OO_F) = |N\cdot\zeta_F(-1)|$. Here $N$ is a
positive integer defined through roots of unity. The formula $(\ast)$ has the factor
$D_F\,2^{[F:\Q]}$ in place of $N$.

\section{Proof}

Write $\OO_\Q$ for the ring of integers of $\Q$, which is identified with $\Z$ by Mathlib's
\texttt{Rat.ringOfIntegersEquiv}.

\begin{lemma}[\texttt{ideal\_eq\_span\_absNorm}, \texttt{absNorm\_span\_nat}, \texttt{card\_ideals\_of\_norm}]
Every ideal $I$ of $\OO_\Q$ equals $(N(I))$, the ideal generated by the natural number $N(I)$, its
absolute norm. For each $m\in\mathbb N$ the ideal $(m)$ has norm $m$. Hence for every $n$ there is
exactly one ideal of norm $n$, so $a_n = 1$.
\end{lemma}

\begin{proof}
$\OO_\Q\cong\Z$ is a principal ideal ring, so $I=(k)$ for an integer $k$. The norm of $(k)$ is
$|N_{\OO_\Q/\Z}(k)| = |k|^{[\Q:\Q]} = |k|$, and $(k) = (|k|)$ since $k=\pm|k|$. The same norm
computation gives $N((m)) = m$. Two ideals of norm $n$ both equal $(n)$.
\end{proof}

\begin{lemma}[\texttt{dedekindZeta\_rat}, \texttt{continuation\_value}, \texttt{riemannZeta\_is\_continuation}]
For $\operatorname{Re}s>1$ the Dirichlet series of $\zeta_\Q$ equals the Riemann zeta function
$\zeta(s)$. Every function $Z$ that is holomorphic on $\C\setminus\{1\}$ and agrees with this series
on $\operatorname{Re}s>1$ satisfies $Z(-1)=\zeta(-1)=-1/12$. The Riemann zeta function itself is such
a $Z$.
\end{lemma}

\begin{proof}
Since $a_n=1$, the series is $\sum_{n\ge1}n^{-s}=\zeta(s)$ (Mathlib's
\texttt{LSeries\_one\_eq\_riemannZeta}). Both $Z$ and $\zeta$ are analytic on the connected open set
$\C\setminus\{1\}$, and they agree on the open half-plane $\operatorname{Re}s>1$. By the identity
theorem they agree on $\C\setminus\{1\}$, which contains $-1$. Finally, Mathlib's
\texttt{riemannZeta\_neg\_nat\_eq\_bernoulli} with $k=1$ gives
$\zeta(-1) = (-1)^1 B_2/2$, and $B_2 = 1/6$ (\texttt{bernoulli\_two}). So $\zeta(-1) = -1/12$
(\texttt{riemannZeta\_neg\_one}).
\end{proof}

\begin{theorem}[\texttt{conjecture\_7928\_false}]\label{thm:main}
Let $Z$ be as in the previous lemma, so $Z(-1)=\zeta_\Q(-1)$. Then
\[
  |Z(-1)|\cdot|d_\Q|\cdot 2^{[\Q:\Q]} = \tfrac{1}{12}\cdot 1\cdot 2 = \tfrac16,
\]
and for every type $G$, $\texttt{Nat.card}\,G \ne 1/6$. In particular $(\ast)$ fails for $F=\Q$,
whatever the group $K_2(\Z)$ is.
\end{theorem}

\begin{proof}
$d_\Q = 1$ (Mathlib's \texttt{Rat.numberField\_discr}) and $[\Q:\Q]=1$
(\texttt{Module.finrank\_self}), so the right-hand side is $1/6$ (\texttt{rhsQ\_riemannZeta}). A
natural number $n$ is either $0$ or at least $1$, so $n\ne 1/6$ (\texttt{nat\_ne\_rhsQ}).
\end{proof}

\section{Lean formalization}

The package \texttt{lean/} (Lean 4, Mathlib \texttt{v4.33.1}) contains
\texttt{Conjecture7928/Basic.lean} (141 lines, namespace \texttt{C7928}). Its only import is
\texttt{import Mathlib}.

\paragraph{Faithfulness of $\zeta_\Q$.} Mathlib's \texttt{NumberField.dedekindZeta K s} is defined
as the raw L-series \texttt{LSeries (fun n => Nat.card \{I : Ideal (O K) // absNorm I = n\}) s}. At
$s=-1$ this series diverges, so its Lean value there is a junk value and not $\zeta_\Q(-1)$. We
therefore do not evaluate it at $-1$. We prove that it equals \texttt{riemannZeta} on
$\operatorname{Re}s>1$. The main theorem then takes \emph{any} function $Z$ holomorphic on
$\C\setminus\{1\}$ that agrees with \texttt{dedekindZeta} $\Q$ on $\operatorname{Re}s>1$, that is, the
continuation, and evaluates $Z(-1)$. The hypotheses are not vacuous: \texttt{riemannZeta} satisfies
them (\texttt{riemannZeta\_is\_continuation}).

\paragraph{Main statements.}
{\small
\begin{verbatim}
theorem conjecture_7928_false (Z : C -> C) (hZ : DifferentiableOn C Z {1}^c)
    (hZeq : forall s : C, 1 < s.re -> Z s = dedekindZeta Q s) (G : Type*) :
    (Nat.card G : R) !=
      ||Z (-1)|| * |((NumberField.discr Q : Z) : R)| * 2 ^ Module.finrank Q Q

theorem conjecture_7928_false_riemannZeta (G : Type*) :
    (Nat.card G : R) !=
      ||riemannZeta (-1)|| * |((NumberField.discr Q : Z) : R)|
        * 2 ^ Module.finrank Q Q

theorem conjecture_7928_false_finite_group (G : Type*) [Group G] [Fintype G] :
    (Fintype.card G : R) != (same right-hand side)
\end{verbatim}
}
(ASCII rendering: \texttt{C}, \texttt{R}, \texttt{Q}, \texttt{Z} stand for $\C,\mathbb R,\Q,\Z$;
\texttt{!=} is $\ne$; \texttt{||.||} is the complex norm.) Supporting lemmas:
\texttt{riemannZeta\_neg\_one}, \texttt{rhsQ\_riemannZeta}, \texttt{ideal\_eq\_span\_absNorm},
\texttt{absNorm\_span\_nat}, \texttt{card\_ideals\_of\_norm}, \texttt{dedekindZeta\_rat},
\texttt{continuation\_value}, \texttt{riemannZeta\_is\_continuation}, \texttt{nat\_ne\_rhsQ}. The
identity-theorem step follows the pattern of Mathlib's own proof in
\texttt{Mathlib/NumberTheory/Harmonic/ZetaAsymp.lean}.

\paragraph{Axiom audit.} \texttt{lake build} succeeds. \texttt{\#print axioms} for each of
\texttt{conjecture\_7928\_false}, \texttt{conjecture\_7928\_false\_riemannZeta},
\texttt{conjecture\_7928\_false\_finite\_group} and \texttt{riemannZeta\_is\_continuation} reports
only \texttt{propext}, \texttt{Classical.choice} and \texttt{Quot.sound}. There is no \texttt{sorry},
no \texttt{native\_decide} and no added axiom.

\section*{Sources}
\begin{itemize}
\item Wikipedia, ``Dedekind zeta function'',
  \url{https://en.wikipedia.org/wiki/Dedekind_zeta_function} (Dirichlet series and meromorphic
  continuation with only a simple pole at $s=1$).
\item Wikipedia, ``Birch--Tate conjecture'',
  \url{https://en.wikipedia.org/wiki/Birch%E2%80%93Tate_conjecture} (context only).
\end{itemize}

\end{document}
